\scanpage{38}

\mspage{- 180 -}\hfill IX

\uline{En d'autres termes}, \uline{la filtration de} $\mathcal{E}^{t}$ \uline{induite par la filtration de Hodge} (\uline{avec les conventions de degrés} (13.4.6)) \uline{est concentrée en degré} $-1$, \uline{celle de} $\mathcal{E}/\mathcal{E}^{f}$ \uline{induite par la filtration de Hodge est concentrée en degré} $0$.

c) \uline{Considérons la filtration de}

\[
\mathcal{E}^{\mathrm{ab}}\overset{\mathrm{dfn}}{=}\mathcal{E}^{f}/\mathcal{E}^{t}\simeq
(\underline{R}^{\mathrm{ab}})\otimes_{\mathbb{Z}}\underline{O}_S
\quad,\quad \text{où }\underline{R}^{\mathrm{ab}}=\underline{R}^{f}/\underline{R}^{t}\quad\text{(cf. a))}\quad,
\]

\uline{induite par la filtration de Hodge de} $\mathcal{E}$,\uline{donc} \uline{définie par le sous-Module} $\underline{F}^{\mathrm{ab}}$ \uline{de} $\mathcal{E}^{\mathrm{ab}}$ \uline{image de}

\[
\underline{F}^{f}\overset{\mathrm{dfn}}{=}\underline{F}\cap\mathcal{E}^{f}\longrightarrow\mathcal{E}^{\mathrm{ab}}\quad ,
\]

(\uline{la flèche étant un monomorphisme grâce à la première relation} (*)).

\uline{Alors le couple} $(\underline{\mathbb{R}}^{\mathrm{ab}},\underline{F}^{\mathrm{ab}}\subset\underline{\mathbb{R}}^{\mathrm{ab}}\otimes_{\mathbb{Z}}\underline{O}_S=\mathcal{E}^{\mathrm{ab}})$ \uline{est isomorphe au couple} $(\underline{R}(B),\underline{F}(B))$ \uline{associé à un} groupe abéloide \uline{sur} $S$, \uline{comme expliqué après} (13.4.4) (\uline{où donc} $B\simeq\mathcal{E}^{\mathrm{ab}}/(F^{\mathrm{ab}}+R^{\mathrm{ab}})$, \uline{les symboles non soulignés désignant les groupes analytiques sur} $S$ \uline{correspondants}).

\uline{Remarque} 13.4.18. Dans la terminologie de DELIGNE [4], la conclusion de la proposition précédente, lorsqu'on se restreint au-dessus d'un ouvert $U$ de $S$ sur lequel $R$ est localement constant, exprime que $\unsure{\mathbb{R}|R}$, muni des deux filtrations (13.4.17.1) de $\mathcal{E}=\underline{\mathbb{R}}\otimes_{\mathbb{Z}}\underline{O}_S$, est une "famille analytique de structures de Hodge mixtes sur $U$ ". La situation qu'on vient d'expliciter a servie de modèle aux conjectures de DELIGNE sur la théorie de Néron en dimensions supérieures, énoncées dans [9, 9.8 à 9.13]. Il y a cependant un trait important de la théorie algébrique

\handwritten{\mspage{498}}

\mspage{498}

% archive sheet number: 37
